\section{Defectos de contrato}
\label{sec:esquema:contrato}\label{sec:evolucion:contrato}

La Tabla~\ref{tab:verificacion} (Sección~\ref{sec:propuesta:prime4api}) recoge lo que ocurre al ingerir las ocho Datasheets en PRIME4API, y sus errores tienen tres causas.
Dos son defectos de contrato, es decir, formas que la guía de ACM2YAML admite y que el motor no procesa.
Son un campo declarado y vacío en Postmark, Resend y Mailchimp Transactional, y un precio que existe y no es público en Resend, SendGrid y Brevo.

La tercera causa es el fallo del propio motor al calcular la curva acumulada de MailerSend y Amazon SES, que mezcla unidades de tiempo (Sección~\ref{sec:propuesta:prime4api}), y no produce propuesta de esquema porque ningún valor de esas Datasheets lo causa.
Mailgun termina con el mismo mensaje, pero su causa no está confirmada.
La forma de \texttt{workload} es un caso distinto, porque ninguna Datasheet del catálogo usa la variante con la que falla el motor.
La guía de ACM2YAML y una Datasheet de ejemplo de su repositorio escriben el campo de dos formas distintas, y se resuelve corrigiendo la guía.

\subsection{La forma de \texttt{workload}: mapa o lista}
\label{subsec:evolucion:dv01}

La guía del esquema documenta \texttt{workload} como un mapa \citep{apisconsumptionmodel}.
Una Datasheet de ejemplo del mismo repositorio lo serializa como una lista, y el TFM de Gavira Sánchez usa también la lista en sus extractos y en el anexo de SendGrid \citep{gavira2026tfm}.
PRIME4API tolera las dos formas en las operaciones de cálculo y solo acepta la lista en las de navegación, que recorren la estructura del documento.
Una Datasheet escrita según la guía produce entonces números correctos y falla cuando la recorren el agente conversacional HARVEY4API o la plataforma SPHERE.

La forma correcta es la lista, porque la exige el modelo de consumo.
Un mapa solo puede declarar una métrica por \textit{endpoint}, y Gavira Sánchez sostiene que un \textit{endpoint} que consume varias métricas a la vez tiene que declararlas por separado \citep{gavira2026tfm}.
Las ocho Datasheets del catálogo usan la lista.
La propuesta es una rectificación: corregir la guía para que documente la lista.
De forma opcional, el motor puede convertir el mapa en lista al leerlo, para que los ficheros escritos con la redacción actual de la guía sigan siendo navegables.

\subsection{P-04. Distinguir «sin límite» de «límite no publicado»}
\label{subsec:evolucion:p04}

\paragraph{El caso.}
Resend no publica el \textit{rate} de su plan Enterprise, y su Datasheet deja presente y vacío el campo que lo declara.
Ese campo vacío dice que el límite existe y que su valor no se conoce.
En Postmark y en Mailchimp Transactional la misma forma dice lo contrario: el proveedor no pone tope en la unidad del campo.
Si solo lee el campo vacío, PRIME4API no puede distinguir el caso de Resend del de Postmark, y P-04 propone cerrar esa ambigüedad.

En Postmark y en Mailchimp Transactional la cota que se publica está en otra métrica, los bytes de la llamada, que es la carencia que cubre P-03 (Apartado~\ref{subsec:evolucion:p02}).
Postmark no pone tope de destinatarios en \texttt{/email/bulk} y lo acota por el tamaño de la petición, 50 MB.
Mailchimp niega por escrito el tope de destinatarios en su API, «there's no recipient limitation», y limita el JSON de cada llamada a menos de 10 MB\footnote{\url{https://mailchimp.com/developer/transactional/docs/outbound-email/}, consultada el 21 de septiembre de 2026.}.
En el sondeo de Mailchimp, una llamada con 5.000 destinatarios devolvió HTTP 200 sin rechazo.

P-04 se promueve por dos vías, y cada una se refiere a un hecho distinto.
La primera es la recurrencia de F-01, un \textit{rate} que existe y no tiene cifra publicada, que cuenta en Postmark, Amazon SES, Mailchimp Transactional y Mailgun (Tabla~\ref{tab:recurrencia}).
En Mailchimp ese \textit{rate} es la cuota horaria de la cuenta de pago, que el proveedor no cifra, y no el tope de destinatarios.
La segunda es la traza: el campo vacío hace fallar al motor en tres Datasheets, aunque solo en Resend es un \textit{rate}.
En Postmark y en Mailchimp es un tope de destinatarios.
Lo que P-04 generaliza es lo que muestra la traza, un campo vacío que el motor no distingue de la ausencia de límite.

\paragraph{La traza.}
PRIME4API convierte a número el valor de esos campos sin comprobar si está vacío, y falla con un error de conversión de tipo.
En Postmark caen por esta causa diez comprobaciones: la consulta por rangos del factor de consumo por petición, tres curvas de capacidad (una es la acumulada), la curva de coste, la recomendación de plan y cuatro operaciones de cálculo.
Esas cuatro dan el tiempo mínimo para alcanzar una capacidad, la capacidad acumulada en un instante, el instante en que se agota la cuota y la espera tras agotarla.
En Resend el mismo error aparece en cinco comprobaciones, y otras ocho fallan por el precio que existe y no es público, el valor \texttt{custom} del plan Enterprise.

Como el fichero de Resend contiene los dos defectos, se generaron dos variantes, cada una con un solo defecto corregido.
Con el campo vacío relleno, la herramienta de verificación de este trabajo da trece errores, todos por el precio no público.
Con el precio no público sustituido por números, da doce, todos por el campo vacío\footnote{Las tres versiones del fichero (la original y las dos variantes) y su salida están en \texttt{03\_datos/\allowbreak apis/\allowbreak resend/\allowbreak evidencias/\allowbreak control\_motor\_causas\_2026-09-23/}.}.
Cada comprobación informa solo del primer error que encuentra.
Por eso, con los dos defectos a la vez, el campo vacío aparece en cinco y el precio no público en ocho, aunque cualquiera de los dos basta por sí solo para hacer fallar doce o trece.

En una tanda aparte de diecisiete llamadas a la API del motor con la Datasheet de Mailchimp Transactional, catorce responden HTTP 500 con el mismo error de conversión, y solo contestan las tres consultas que no calculan tiempo.
Si el campo vacío del tope se rellena en el documento ya cargado por el motor, sin tocar el fichero, el error pasa a ser una división por cero.
Esa segunda caída la causa otro campo vacío, el periodo de la franquicia de la cuenta de prueba, 500 correos gratuitos que no caducan ni se renuevan (F-33).
Con los dos campos rellenos, la misma llamada devuelve HTTP 200.
La herramienta de verificación de este trabajo, con sus propias comprobaciones, da a Mailchimp diez errores con el mismo mensaje que a Postmark.

\paragraph{Por qué el esquema no ofrece salida.}
ACM2YAML v0.3 deja dos alternativas, y las dos tienen coste.
Omitir el campo hace que la Datasheet afirme implícitamente que no hay límite.
En Resend, donde el límite existe, la omisión produciría una respuesta optimista y falsa.
En Postmark y en Mailchimp coincidiría con lo que dice el proveedor, pero PRIME4API no distinguiría esa omisión de la que haría Resend con un límite que existe.
Declarar el campo vacío, la otra alternativa, deja sin responder operaciones enteras del motor: diez comprobaciones en Postmark, cinco en Resend y catorce de diecisiete llamadas en Mailchimp.

La Datasheet de Amazon SES elige la omisión para el caudal de una cuenta con acceso de producción: no lo declara y anota la ausencia en su registro \texttt{not\_found}.
Ese registro, que usan las ocho Datasheets, dice qué falta en el documento, pero no en qué campo, y una marca junto al propio campo dejaría esa anotación donde el motor la puede leer.

\paragraph{Propuesta.}
Se propone una marca explícita de valor no publicado junto al campo.
El esquema declara la marca y el motor la respeta: omite esa dimensión y sigue evaluando el resto en lugar de abandonar el documento.
Con la marca, el límite sin cifra de Resend y la ausencia de tope de Postmark y Mailchimp dejan de escribirse igual.
Resend la lleva en su caudal de Enterprise.
Postmark y Mailchimp pueden omitir el campo, y la omisión dice lo mismo que su fuente: no hay tope de destinatarios.
La cota en bytes ya la declara la Datasheet de Mailchimp como un \texttt{workload} aparte, y la de Postmark necesita P-03.

La misma distinción alcanza al periodo de una cuota, el segundo campo vacío de Mailchimp.
Ahí, igual que en su tope de destinatarios, no falta una cifra, porque la franquicia de 500 correos de la cuenta de prueba no caduca ni se renueva (Sección~\ref{sec:propuesta:catalogo}).
La marca queda para los valores que existen y no se publican, y la cuota sin periodo es otro elemento, F-33.

\subsection{P-05. Declarar un plan cuyo precio existe y no es público}
\label{subsec:evolucion:p05}

\paragraph{El caso.}
Resend no publica el precio ni la cuota de su plan Enterprise, y la Datasheet escribe en los dos campos el valor \texttt{custom}.
Es lo que publica el proveedor, y es la forma que la guía documenta para un valor negociable \citep{apisconsumptionmodel}.

\paragraph{La traza.}
PRIME4API intenta convertir la cadena \texttt{custom} a número y falla.
El fallo no se queda en el plan Enterprise.
La evaluación recorre todos los planes del documento, así que ese único valor impide evaluar también los planes Free, Pro y Scale, que tienen precio publicado.
Con ese error fallan las tres consultas de \textit{rates}, cuotas y límites, las cuatro operaciones de cálculo y la consulta del precio del plan Enterprise, que son ocho de los trece errores de la Datasheet de Resend.
Con el campo vacío relleno, los trece pasan a esta causa (Apartado~\ref{subsec:evolucion:p04}).
SendGrid y Brevo reproducen el fallo con trece errores cada una, todos con el mismo mensaje: la cadena \texttt{custom} no se puede convertir a número.

\paragraph{Por qué importa.}
En siete de las ocho APIs del catálogo hay al menos un precio que el proveedor no publica, y la excepción es Amazon SES.
En Postmark, Resend, SendGrid, MailerSend, Brevo y Mailgun es un plan a medida, y en Mailchimp Transactional es el precio por encima de cuatro millones de correos al mes.
Resend, SendGrid y Brevo escriben ese plan con \texttt{custom}, como admite la guía, y el motor falla en las tres.
Postmark, MailerSend y Mailgun lo dejan fuera de sus planes.

\paragraph{Propuesta.}
Se proponen dos cambios separables.
En el esquema, una forma explícita de declarar un plan que tiene precio aunque no lo publique, en sustitución de la cadena \texttt{custom} que el motor no sabe leer.
En el motor, aislar el fallo por plan, de modo que un plan no evaluable devuelva su error y los demás se sigan calculando.
El segundo cambio no depende del primero, y por sí solo recupera tres de los cuatro planes de la Datasheet de Resend.

Las dos propuestas de esta sección se han medido solo contra PRIME4API.
Otro motor podría tolerar el campo vacío o la cadena \texttt{custom} sin que el esquema cambiara.
